\begin{recipe}{Zelfgemaakte vla}
\kicker[Dessert,Zoet,Basisrecept]{Dessert · zoet · basisrecept}
\index[register]{Zelfgemaakte vla}
\index[register]{Vanille}
\index[register]{Eidooiers}
\index[register]{Maïzena}

\meta{15 minuten \dvd Ongeveer 1 liter \dvd Eenvoudig}

\lettrine{W}{e} hebben deze vla ooit gemaakt puur om te zien hoe dat nou eigenlijk
gaat, en het bleek best makkelijk. Een paar basisingrediënten is al genoeg. Reken
wel op een paar uur afkoeltijd voor je een romige vanillevla op tafel hebt staan.

\blockrule{Ingrediënten}
\begin{ingredients}
  \ing{1 liter}{volle melk}
  \ing{75 g}{suiker}
  \ing{40 g}{maïzena}\index[register]{Maïzena}
  \ing{2}{eidooiers}\index[register]{Eidooiers}
  \ing{1}{vanillestokje}\index[register]{Vanille}
  \ingb{snufje zout}
\end{ingredients}

\blockrule{Bereiding}
\tip{Leg de plasticfolie echt rechtstreeks op de vla, zonder lucht ertussen. Zo
voorkom je een taai velletje.}
\begin{steps}
  \item Giet 900 ml van de melk in een pan en houd de resterende 100 ml apart.
  \item Snijd het vanillestokje overlangs open en schraap het merg eruit. Voeg
        het merg en het lege stokje toe aan de melk in de pan. Breng de melk
        op middelhoog vuur langzaam tegen de kook aan en laat op laag vuur
        10 minuten trekken.
  \item Meng in een kom de suiker, de maïzena en het snufje zout. Roer de
        apart gehouden 100 ml koude melk erdoor tot een glad papje en klop er
        de eidooiers goed doorheen.
  \item Haal het vanillestokje uit de warme melk. Giet al roerend een flinke
        scheut van de warme melk bij het eimengsel en roer direct goed door,
        zodat de eieren niet gaan stollen.
  \item Giet het verdunde eimengsel terug in de pan bij de rest van de warme
        melk. Breng het geheel op laag vuur onder voortdurend roeren met een
        garde zachtjes aan de kook en laat het 1 tot 2 minuten doorkoken tot
        de vla dikker wordt en bindt.
  \item Giet de warme vla door een fijne zeef in een schaal. Dek de vla direct
        af met een stuk plasticfolie dat je rechtstreeks op de vla legt, zodat
        er geen vel ontstaat. Laat volledig afkoelen in de koelkast, minimaal
        2 uur.
  \item Klop de vla voor het serveren nog kort door met een garde voor een
        extra glad resultaat.
\end{steps}
\end{recipe}
